\section{What survives for SGD and resource-bounded learners}\label{sec:resources}
The SQ theorem does not settle the benign-SGD question. A generic SQ algorithm can store or regenerate its class description for free before issuing any query. Ordinary gradient dynamics have a fixed network, initialization, and objective. The SQ-to-differentiable-learning simulation of \citet{AKMSS21} uses engineered models and initialization; its computational network size depends on the simulated computation, not solely on the number of SQ calls. Its ReLU alternative also allows fixed weights. The table-polynomial implementation may take exponential time in $n$. The succinct implementation has polynomial time in $n$, but its key length and runtime exponents grow with the chosen sign-rank power. Neither implementation by itself proves a smaller-than-dimension plain network with the stipulated initialization and dynamics.

For reference, a robust population-gradient method with $S$ parameters, $T$ calls, per-example coordinate bound $B_g$, and adversarial Euclidean error tolerance $\delta$ is simulated by
\[
 m=TS,\qquad\tau=\delta/(B_g\sqrt S).
\]
Query each gradient coordinate divided by $B_g$, then multiply back: the coordinate errors are at most $\delta/\sqrt S$, so their squared sum is at most $\delta^2$. Induction over updates proves the simulation. Conversely the adaptive objective $J_q(t;x,y)=tq(x,y)$ has population derivative $\E q$, realizing a scalar SQ. This converse permits changing the objective; it is not the fixed logistic-loss process of the SGD question.

The remaining SGD question is exactly the one with $X=\{-1,+1\}^n$, a fully connected bias-free ReLU network, independent layer-$i$ Gaussian initialization of variance $1/n_{i-1}$, online gradients of $\log(1+e^{-yf_\theta(x)})$, one fixed step size, and classification by the specified tail sum. Does a common expected-error guarantee at $\epsilon<1/4$ imply $\dc(H)\le CTS$ or $\dc(H)\le\operatorname{poly}(T,S)$? A sampled gradient is not a worst-case approximate population gradient, and success under its sampling law does not imply robustness to every vector perturbation. No answer for those particular dynamics is established here.

\begin{samepage}
\begin{proposition}[Representation does not give the required dynamics]\label{thm:network}
Pad the grid domain to a cube of dimension $n=\lceil\log_2(2N^3)\rceil$, giving every padding point label $+1$. Every member is represented exactly by a bias-free ReLU network with widths $(n,2n,N+1,1)$ and at most
\[
 S=2n^2+2n(N+1)+(N+1)
\]
weights. Its common dimension complexity is at most $2N+1$. This proposition is an upper bound, not a claim that $S$ parameters are necessary.
\end{proposition}
\end{samepage}
\begin{proof}
The first hidden layer computes $\sigma(x_i),\sigma(-x_i)$ for every coordinate. These give $x_i$ by subtraction and the constant $c=\sigma(x_1)+\sigma(-x_1)=1$ on the cube. For every negative point $p$, a second-layer unit computes
\[
 \sigma\left(\sum_i p_ix_i-(n-1)c\right).
\]
At $x=p$ its value is one; at every other cube point the sum is at most $n-2$ and the unit is zero. Pass $c$ through one additional unit and output $c$ minus twice the sum of the negative-point units. Pad unused units by zero weights. There are at most $N$ negative points, and counting all fully connected weights gives the claimed $S$. The moment-curve proof in Theorem~\ref{thm:construction} still applies to the padded domain. Padding preserves every dimension lower bound by restriction; it preserves a rectangle lower bound up to a factor two, because either padding carries at least half the point mass or the original domain does.
\end{proof}

A typical flip table has $I=\Theta(N^4)$ unconstrained bits. This does not by itself prove a lower bound of $N$ on the number of unrestricted real network parameters: a real parameter has no fixed bit budget. Under a separate $b$-bit-per-weight convention, a fixed $S$-weight architecture has at most $2^{bS}$ parameter choices, so representing all $2^N$ independently specified patterns on a fixed line would require $bS\ge N$. That bounded-precision counting statement concerns an entire family of patterns, not a necessary parameter count for every fixed row of one fixed $A$. It is not used to answer the SGD question.

\begin{samepage}
\begin{theorem}[Source-description length alone does not repair the implication]\label{thm:description}
Interpret $\ell$ as the bit length of a finite learner program, including literal hard-coded constants and data, but do not charge running time or generated workspace. There is a computable sequence of incidence-flip classes and learners with
\[
 n,\ell,m=O(\log N),\qquad 1/\tau=O_\epsilon(1),\qquad
 \dc(H)=\Omega(N/\log N).
\]
Consequently no universal bound $\dc(H)\le\operatorname{poly}(m,1/\tau,n,\ell)$ holds in this interpretation.
\end{theorem}
\end{samepage}
\begin{proof}
Let $K=2N^3$, $b_N=\lceil\log_2(4K)\rceil$, and $d_N=\lfloor N/(128b_N)\rfloor$. For sufficiently large $N$, the counting bound in Theorem~\ref{thm:construction} guarantees a flip matrix of sign-rank greater than $d_N$: indeed $\log_2(4eK)\le2b_N$, so
\[
 2Kd_N\log_2(4eK)\le N^4/16<I.
\]
Select the lexicographically first such matrix. This selection is computable. For a proposed matrix $A$, the predicate $\sr(A)\le d_N$ is the truth of the finite real sentence
\[
 \exists U\in\R^{K\times d_N},\ V\in\R^{d_N\times K}\quad
 \bigwedge_{i,j} A_{ij}\sum_{s=1}^{d_N}U_{is}V_{sj}>0.
\]
Decision of finite polynomial real sentences is the classical real-closed-field decision theorem. We use this explicit imported decidability result, not a claimed practical low-rank solver; a formal constructive treatment is given by \citet{MC12}. Enumerate the finitely many flip tables and decide the displayed sentence for each until the first failure. The counting argument ensures termination. Small $N$ can be handled by any fixed table.

A single fixed program performs this generation and then runs the finite rational rectangle-mixture algorithm. Encode $N$ in $O(\log N)$ bits. Fix an accuracy such as $1/16$, so that all accuracy constants have constant-length descriptions. Conservative rational ceilings can replace the logarithmic constants in Theorem~\ref{thm:main}, making every step an ordinary finite computation. The table generation happens before the first SQ; it has no query cost. The learner therefore has $m=O(\log K)$ and constant positive tolerance, while the chosen class has $\dc>d_N=\Omega(N/\log N)$. Its program description has size $O(\log N)$. The generated table occupies far more space, which this interpretation does not charge. The displayed polynomial implication is contradicted.
\end{proof}

This does not produce an efficient explicit counterexample. It instead shows that a resource-aware repair must charge computation, generated memory, or an expanded circuit description, not just the length of the program that generates them. One genuinely stronger question is whether charging total preprocessing and query-evaluation time $t$ yields a bound polynomial in $(m,1/\tau,n,\ell,t)$ in a specified finite-bit model. We state this as an open question, not as an established or uniquely natural repair.

\begin{samepage}
\begin{proposition}[What the classical-method barrier actually bounds]\label{thm:classicalbarrier}
For the orientation in which rows are hypotheses, let
\[
 m_{\rm avg}(A)=\inf_{\mu\in\Delta(X)}\ 
 \sup_{\substack{\text{strict unit-vector}\text{representations of }A}}
 \min_h\E_{x\sim\mu}|\langle v_h,u_x\rangle|.
\]
Then
\[
 \sqrt{\operatorname{VC}(H)}\le m_{\rm avg}(A)^{-1}\le\rect(A)^{-1}.
\]
In particular, if $\rect(A)^{-1}\le\operatorname{poly}(n)$, these three classical lower-bound parameters cannot certify superpolynomial sign-rank. This does not upper-bound sign-rank itself and does not imply that every SQ-easy class has large rectangles.
\end{proposition}
\end{samepage}
\begin{proof}
Fix a point prior $\mu$. Apply the finite separation argument of Lemma~\ref{lem:mixture} to the transposed matrix, obtaining a distribution $\lambda$ over monochromatic rectangles $S_R\times T_R$ with
\[
 \forall h,\quad\E_R[\1_{T_R}(h)\mu(S_R)]\ge r,
 \qquad r=\rect(A).
\]
In the Euclidean space indexed by rectangles, put
\[
 a_h(R)=c_R\sqrt{\lambda_R}\1_{T_R}(h),\qquad
 b_x(R)=\sqrt{\lambda_R}\1_{S_R}(x).
\]
Their norms are at most one, and their inner product is $A_{hx}$ times the probability that $R$ covers $(h,x)$. These products may be zero, so this alone is not a strict representation. Add an orthogonal strict unit-vector representation scaled by $\sqrt\delta$, and scale the displayed vectors by $\sqrt{1-\delta}$. Every resulting product is strict with the required sign, and its average absolute value in each row is at least $(1-\delta)r$. Pad row norms and point norms separately in orthogonal private coordinates to make them one without changing any cross inner products. Letting $\delta\downarrow0$ proves $m_{\rm avg}(A)\ge r$.

For the VC bound, suppose $v$ points are shattered and use their uniform point prior. In any strict unit-vector representation, select one row $h_\sigma$ for each sign pattern $\sigma\in\{-1,+1\}^v$. If the minimum row-average margin is $a$, then
\[
 va\le\sum_{i=1}^v\sigma_i\langle v_{h_\sigma},u_i\rangle
 \le\left\|\sum_{i=1}^v\sigma_i u_i\right\|.
\]
Averaging the square of the last norm over independent fair signs gives $v$, because cross terms cancel. Some sign pattern has norm at most $\sqrt v$, so $a\le1/\sqrt v$. Take the supremum over representations and the infimum over priors. For $v=0$ the inequality is immediate. This proves the chain, including the strict-sign correction.
\end{proof}

The chain and its original average-margin interpretation are due to \citet{HHPTZ22}; the proof is included to fix the orientation and strict-sign issue. A counterexample obtained with a polynomial inverse rectangle bound cannot be certified as high-sign-rank by those numerical parameters. Counting is already a way around that barrier, as Theorem~\ref{thm:construction} demonstrates. Therefore the barrier does not rescue the program-length conjecture. For an \emph{efficiently explicit} sequence it motivates seeking a method beyond those parameters, in the direction of \citet[Problem~4.1]{HHPTZ22}.

The recent topological methods provide a distinct parameter, not one proved bounded by this chain. The checked results do not supply the missing efficient total construction: \citet{FHV26} prove near-$2k$ bounds for certain \emph{partial} Gap Hamming matrices, while \citet{GHIS25} bound the sign-rank of total $k$-Hamming-distance matrices by $2^{O(k)}$ independently of the bit dimension. These facts neither rule out a future topological route nor prove that it evades the rectangle obstruction on the particular flip family. No such explicit corollary is claimed here.

\subsection{Unknown tables and hidden keys}\label{sec:secret}
The following bounds quantify the cost of hiding the class. Their quantifier is one learner chosen independently of the flip table or key. That is different from choosing an architecture separately for each fixed class.

\begin{samepage}
\begin{proposition}[Unseen independent labels]\label{prop:unseen}
Fix the full-length template line $\ell_0:y=x+1$ and the uniform distribution on its $N$ points. Choose $A$ by independent fair incidence flips. Let a single learner, chosen independently of $A$, receive $T$ independent labeled samples from $(D,h_{\ell_0})$ and no further information about $A$. Then
\[
 \E_{A,\mathrm{sample},\mathrm{coins}}L_{D,h_{\ell_0}}(\widehat h)
 \ge\frac12(1-1/N)^T
 \ge\frac12(1-T/N).
\]
Thus a guarantee of average error at most $\epsilon$, or a guarantee for every flip table, requires $T\ge(1-2\epsilon)N$. If $\epsilon<1/4$ and $S\ge2$, such a uniformly successful $S$-parameter learner satisfies $TS\ge\dc(H_A)/3$ for every member of the flip family.
\end{proposition}
\end{samepage}
\begin{proof}
Condition on the sampled point indices, their observed labels, and the learner's random coins. The label at any unobserved point of $\ell_0$ remains an independent fair bit. The learner's prediction there has conditional error $1/2$. A fixed point is unobserved with probability $(1-1/N)^T$. Average over the $N$ points to get the first inequality. The second follows by induction from $(1-z)^T\ge1-Tz$ for $0\le z\le1$. Rearranging proves the sample requirement. For the final claim, $T>N/2$ and $S\ge2$ give $TS\ge N$, whereas $\dc(H_A)\le2N+1\le3N$ for every table.
\end{proof}

The lower bound is an average over tables, or a minimax statement against a learner that must work for all of them. It is false as a requirement for every individual fixed table. For example, the all-positive table needs no sample. The argument remains valid even if the learner is told $\ell_0$ and $D$: the unavailable information is the unobserved independent labels, not the line identity.

\begin{samepage}
\begin{proposition}[Hidden-key computational hardness]\label{prop:hiddenkey}
Under the PRF assumption of Theorem~\ref{thm:prf}, no probabilistic polynomial-time learner independent of $s$, with a polynomial-time evaluable output and only polynomially many labeled samples from the uniform distribution on $\ell_0$, has average error at most $1/2-\gamma(n)$ for uniform $s$, where $\gamma(n)\ge1/\operatorname{poly}(n)$. More precisely, a learner with $q(n)$ samples and advantage $\gamma(n)$ gives a PRF distinguisher of advantage at least
\[
 \gamma(n)-\frac{q(n)}{2N}.
\]
\end{proposition}
\end{samepage}
\begin{proof}
Given an oracle function, generate the learner's independent uniform point samples on $\ell_0$ and label them by its oracle values at the corresponding incidence encodings. Run the learner without giving it a key. Draw an independent uniform test point on $\ell_0$, query its label, and output one if the predictor agrees. In the PRF case acceptance is at least $1/2+\gamma(n)$. In the random-function case, conditional on the training transcript, a previously unqueried point has an independent fair label. A fresh test point repeats a training point with probability at most $q(n)/N$. Acceptance is therefore at most $1/2+q(n)/(2N)$. The simulation and one output evaluation are polynomial-time. Since $N=2^{(n-1)/3}$, the collision term is negligible, contradicting PRF security for inverse-polynomial $\gamma$.
\end{proof}

A guarantee for every key, or for all but a negligible fraction of keys, would imply the prohibited average guarantee. The proposition does not exclude an arbitrary algorithm specialized to one particular key. Nor does it automatically apply to infinite-precision real computation, which is not the finite-bit model in the cryptographic assumption.

\paragraph{The exact remaining SGD question.}
The original SGD premise permits the architecture, step size and horizon to be chosen for each class $H_A$. Consequently it does not itself impose independence from $A$, even though the weights are initialized from the prescribed Gaussian law. The two preceding propositions do not reverse those quantifiers. They prove obstructions for a \emph{single key-independent learner}, not an unconditional lower bound for every class-dependent benign-SGD design. Our SQ algorithms use the class through a table or an evaluable key; their generic differentiable simulations do not preserve the bias-free fully connected Gaussian-initialized logistic-SGD process. A negative answer for that process requires a compatible efficient representation and a proof of its dynamics for every marginal, not merely an oracle learner or a hidden-key lower bound. For an efficiently explicit class having a large rectangle ratio, Proposition~\ref{thm:classicalbarrier} also explains why the listed classical certificate parameters cannot establish superpolynomial sign-rank; it does not impose that barrier on every possible SQ-easy or SGD-easy class. Thus the original class-dependent benign-SGD question, and its stronger uniform key-independent variant, are stated separately rather than declaring the negative direction unreachable.

